% GENERATED FILE. Edit canonical lessons, then run npm run generate:books.
% canonical-source-hash: fnv1a64:57e9b3ec24e47553
% canonical-lessons: TE-C169-pundu, TE-C169-bobba, TE-C169-daddurlu, TE-C169-macca, TE-C169-benuku

\chapter{Sore, Blister, Rash, Scar, Sprain}
\label{ch:te-sore-blister-rash-scar-sprain}
\hlchaptermodality{169}

\begin{chapteropening}
\textbf{By the end of this chapter:} \emph{I can say a sore, a blister, a rash, a scar and a sprain.}

Complete the last lesson of chapter 169: I can say a sore, a blister, a rash, a scar and a sprain.
\end{chapteropening}

\section[\texorpdfstring{\te{పుండు} (puṇḍu)}{పుండు (puṇḍu)}]{\te{పుండు} (puṇḍu) — a sore, an ulcer}

\label{lesson:TE-C169-pundu}



\begin{quote}
Before the new one: say the Telugu for peace, then the Telugu for patience.
\end{quote}

\subsection*{You'll want to know: \te{పుండు}}
\textbf{\te{పుండు}} — \emph{puṇḍu} — \textquotedblleft{}a sore, an ulcer\textquotedblright{}.

\begin{grammarlens}[title={What you've built}]
One more word for this chapter.
\end{grammarlens}

\subsection*{Your turn}
\begin{itemize}
\raggedright
  \item \emph{Say it:} \emph{puṇḍu}
  \item \emph{Say it:} \emph{puṇḍu}, once more
  \item \emph{Say it:} say \emph{puṇḍu}
  \item \emph{Recall:} say the Telugu for peace, then the Telugu for patience, then say \emph{puṇḍu} again
\end{itemize}

\begin{tcolorbox}[breakable,colback=teal!4,colframe=teal!35!black,title={Before you move on}]
What does \te{పుండు} mean? (\textquotedblleft{}A sore, an ulcer\textquotedblright{}.) Say it once more.
\end{tcolorbox}
\section[\texorpdfstring{\te{బొబ్బ} (bobba)}{బొబ్బ (bobba)}]{\te{బొబ్బ} (bobba) — a blister}

\label{lesson:TE-C169-bobba}



\begin{quote}
Before the new one: say the Telugu for patience, then the Telugu for a sore.
\end{quote}

\subsection*{You'll want to know: \te{బొబ్బ}}
\textbf{\te{బొబ్బ}} — \emph{bobba} — \textquotedblleft{}a blister\textquotedblright{}.

\begin{grammarlens}[title={What you've built}]
One more word for this chapter.
\end{grammarlens}

\subsection*{Your turn}
\begin{itemize}
\raggedright
  \item \emph{Say it:} \emph{bobba}
  \item \emph{Say it:} \emph{bobba}, once more
  \item \emph{Say it:} say \emph{bobba}
  \item \emph{Recall:} say the Telugu for patience, then the Telugu for a sore, then say \emph{bobba} again
\end{itemize}

\begin{tcolorbox}[breakable,colback=teal!4,colframe=teal!35!black,title={Before you move on}]
What does \te{బొబ్బ} mean? (\textquotedblleft{}A blister\textquotedblright{}.) Say it once more.
\end{tcolorbox}
\section[\texorpdfstring{\te{దద్దుర్లు} (daddurlu)}{దద్దుర్లు (daddurlu)}]{\te{దద్దుర్లు} (daddurlu) — a rash}

\label{lesson:TE-C169-daddurlu}



\begin{quote}
Before the new one: say the Telugu for a sore, then the Telugu for a blister.
\end{quote}

\subsection*{You'll want to know: \te{దద్దుర్లు}}
\textbf{\te{దద్దుర్లు}} — \emph{daddurlu} — \textquotedblleft{}a rash\textquotedblright{}.

\begin{grammarlens}[title={What you've built}]
One more word for this chapter.
\end{grammarlens}

\subsection*{Your turn}
\begin{itemize}
\raggedright
  \item \emph{Say it:} \emph{daddurlu}
  \item \emph{Say it:} \emph{daddurlu}, once more
  \item \emph{Say it:} say \emph{daddurlu}
  \item \emph{Recall:} say the Telugu for a sore, then the Telugu for a blister, then say \emph{daddurlu} again
\end{itemize}

\begin{tcolorbox}[breakable,colback=teal!4,colframe=teal!35!black,title={Before you move on}]
What does \te{దద్దుర్లు} mean? (\textquotedblleft{}A rash\textquotedblright{}.) Say it once more.
\end{tcolorbox}
\section[\texorpdfstring{\te{మచ్చ} (macca)}{మచ్చ (macca)}]{\te{మచ్చ} (macca) — a scar, a spot}

\label{lesson:TE-C169-macca}



\begin{quote}
Before the new one: say the Telugu for a blister, then the Telugu for a rash.
\end{quote}

\subsection*{You'll want to know: \te{మచ్చ}}
\textbf{\te{మచ్చ}} — \emph{macca} — \textquotedblleft{}a scar, a spot\textquotedblright{}.

\begin{grammarlens}[title={What you've built}]
One more word for this chapter.
\end{grammarlens}

\subsection*{Your turn}
\begin{itemize}
\raggedright
  \item \emph{Say it:} \emph{macca}
  \item \emph{Say it:} \emph{macca}, once more
  \item \emph{Say it:} say \emph{macca}
  \item \emph{Recall:} say the Telugu for a blister, then the Telugu for a rash, then say \emph{macca} again
\end{itemize}

\begin{tcolorbox}[breakable,colback=teal!4,colframe=teal!35!black,title={Before you move on}]
What does \te{మచ్చ} mean? (\textquotedblleft{}A scar, a spot\textquotedblright{}.) Say it once more.
\end{tcolorbox}
\section[\texorpdfstring{\te{బెణుకు} (beṇuku)}{బెణుకు (beṇuku)}]{\te{బెణుకు} (beṇuku) — a sprain}

\label{lesson:TE-C169-benuku}



\begin{quote}
Before the new one: say the Telugu for a rash, then the Telugu for a scar.
\end{quote}

\subsection*{You'll want to know: \te{బెణుకు}}
\textbf{\te{బెణుకు}} — \emph{beṇuku} — \textquotedblleft{}a sprain\textquotedblright{}.

\begin{grammarlens}[title={What you've built}]
One more word for this chapter.
\end{grammarlens}

\subsection*{Your turn}
\begin{itemize}
\raggedright
  \item \emph{Say it:} \emph{beṇuku}
  \item \emph{Say it:} \emph{beṇuku}, once more
  \item \emph{Say it:} say \emph{beṇuku}
  \item \emph{Recall:} say the Telugu for a rash, then the Telugu for a scar, then say \emph{beṇuku} again
\end{itemize}

\begin{tcolorbox}[breakable,colback=teal!4,colframe=teal!35!black,title={Before you move on}]
What does \te{బెణుకు} mean? (\textquotedblleft{}A sprain\textquotedblright{}.) Say it once more.
\end{tcolorbox}
